\documentclass[11pt,a4paper]{article}
\usepackage{fontspec}
\IfFontExistsTF{Times New Roman}{\setmainfont{Times New Roman}}{\setmainfont{DejaVu Serif}}
\usepackage[russian]{babel}
\usepackage{amsmath,amssymb,amsthm,booktabs}
\usepackage[margin=24mm,top=21mm,bottom=22mm]{geometry}
\usepackage{xcolor}
\definecolor{ink}{HTML}{1B4152}
\usepackage[unicode,colorlinks=true,urlcolor=ink,linkcolor=ink,citecolor=ink]{hyperref}
\hypersetup{pdftitle={Третья точка наблюдения в цепочечной камерной модели},pdfauthor={Mikita Aharodni}}
\newtheorem{thm}{Теорема}
\newcommand{\Res}{\operatorname{Res}}
\newcommand{\diag}{\operatorname{diag}}
\newcommand{\dd}{\mathrm d}
\newcommand{\coef}[2]{[#1]\,#2}
\renewcommand{\proofname}{Доказательство}
\setlength{\emergencystretch}{2em}
\raggedbottom
\begin{document}
\begin{center}
{\LARGE\bfseries\color{ink} Третья точка наблюдения\\ в цепочечной камерной модели\par}
\vspace{7pt}Mikita Aharodni\\[2pt]
{\small Рабочая исследовательская записка, 26 ноября 2026 г.}
\end{center}
\begin{abstract}
В цепочке с входом в первую камеру наблюдения в третьей и последней камерах могут допускать два положительных регулярных набора скоростей. Доказывается, что добавление любой третьей камеры даёт глобальную единственность для типичных положительных параметров. Измерение в камере 1, 2 или 4 даёт единственность во всём положительном ортанте, когда измеряемые камеры различны. Для более удалённых выходов исключения действительно встречаются. Приводится бесконечное семейство и воспроизводимый семикамерный пример с различными скоростями выведения.
\end{abstract}

\section{Постановка и результат}
Пусть \(n\ge4\). В модели \(x'=Ax+e_1u,\ x(0)=0\) скорости \(a_j,b_j>0\) относятся к переходам \(j\to j+1\) и \(j+1\to j\). Единственное выведение скорости \(h>0\) происходит в известной камере \(\ell\). Поэтому \(A_{j+1,j}=a_j,\ A_{j,j+1}=b_j,\ A_{jj}=-d_j\), где \(d_j=a_j+b_{j-1}+h\mathbf1_{j=\ell}\), \(a_n=b_0=0\). Наблюдаются полные точные кривые количеств вещества при известном входе. Равенство данных означает равенство трёх передаточных функций.

\begin{thm}\label{t:main}
Пусть исходные выходы расположены в камерах \(3,n\), а дополнительный выход --- в \(r\notin\{3,n\}\).
\begin{enumerate}
\item Для всякого \(r\) вне собственного алгебраического множества положительных исходных параметров существует единственный положительный вектор скоростей, дающий те же три выхода. Конкурирующий вектор может лежать и в исключительном множестве.
\item Если требуется гарантия при всех положительных параметрах и при любом известном месте выведения, при \(4\le n\le6\) подходит каждый допустимый \(r\). При \(n\ge7\) подходят в точности \(r=1,2,4\).
\item Для каждой пары \(n\ge7,\ 5\le r<n\) исключение реализуется уже при \(\ell=1\): существуют два различных положительных вектора, регулярных для исходной карты двух выходов, с одинаковыми тремя выходами.
\end{enumerate}
\end{thm}

Положим \(\beta_j=a_jb_j\). Матрица \(-A\) диагонально подобна симметричной \(T\) с диагональю \(d_j\) и внедиагональными элементами \(-\sqrt{\beta_j}\).
Через \(L_j,R_j\) обозначим полиномы \(\det(sI+T)\) для интервалов \([1,j]\) и \([j,n]\); пустой интервал имеет определитель 1. Тогда
\[
 D=L_n,\quad W=R_4,\quad g_j=\prod_{v<j}a_v,\qquad
 G_j(s)=\frac{g_jR_{j+1}(s)}{D(s)},\qquad p=g_3,\quad q=g_n .
\]
Последнее равенство для \(G_j\) следует из разложения кофактора по пути от 1 до \(j\). Числитель \(G_n\) --- положительная константа \(q\), так что \(D\) восстанавливается целиком. В \(D G_j\) старший коэффициент даёт \(g_j\); деление на него даёт монический \(R_{j+1}\), в том числе после сокращения исходной дроби.

\newpage
\section{Точки, достаточные при всех параметрах}
Соседние хвостовые полиномы взаимно просты. Действительно,
\(R_j=(s+d_j)R_{j+1}-\beta_jR_{j+2}\), \(\beta_j>0,\ R_{n+1}=1\). Общий корень распространялся бы до константы. Деление двух соседних полиномов даёт \(d_j\), затем \(\beta_j\) как старший коэффициент остатка с противоположным знаком, и \(R_{j+2}\). Аналогичное правило слева восстанавливает коэффициенты по \(L_j,L_{j-1}\).

Если \(r=1\), известны \(D=R_1,R_2\); последовательное деление восстанавливает весь \(T\). Пусть среди выходов есть соседние \(k,k+1\), так что известны \(R_{k+1},R_{k+2}\). Разложение по ребру \(k\) даёт
\begin{equation}\label{eq:split}
D=L_kR_{k+1}-\beta_k L_{k-1}R_{k+2}.
\end{equation}
Из двух решений следовало бы \(R_{k+1}\mid\Delta(\beta_k L_{k-1})\). При \(n\ge2k\) степень делителя \(n-k\) больше \(k-1\). Поэтому совпадают \(\beta_k L_{k-1}\) и \(L_k\). Моничность \(L_{k-1}\) определяет \(\beta_k\); деление полиномов с двух концов даёт все \(d,\beta\). Для \(r=2\) берём \(k=2\), что покрывает все \(n\ge4\); для \(r=4\) берём \(k=3\), что покрывает \(n\ge6\).

Остаётся \(n=5,\ r=4\). Из \(R_4,R_5\) известно \(d_4\); из \(g_3,g_4,g_5\) известны \(a_3=g_4/g_3,\ a_4=g_5/g_4\). При \(\ell=4\) баланс массы при постоянном единичном входе даёт \(h=1/G_4(0)\). Поэтому при всяком \(\ell\) известна
\[
 \beta_3=a_3(d_4-a_4-h\mathbf1_{\ell=4}).
\]
Теперь \eqref{eq:split} даёт \(R_4\mid\Delta L_2\); моничность понижает степень разности до единицы. Следовательно, \(L_2,L_3\) однозначны. Баланс в нуле допустим, поскольку
\[
 v^\top Tv=\sum_{j<n}(\sqrt{a_j}v_j-\sqrt{b_j}v_{j+1})^2+h v_\ell^2>0
 \quad(v\ne0).
\]

При известной \(\ell\) переход от \(T\) к скоростям однозначен. Слева от \(\ell\) вычисляем \(a_j=d_j-b_{j-1}\), затем \(b_j=\beta_j/a_j\). Справа, в обратном порядке, вычисляем \(b_{j-1}=d_j-a_j\), затем \(a_{j-1}=\beta_{j-1}/b_{j-1}\). Наконец \(h=d_\ell-a_\ell-b_{\ell-1}\). Во всех этих выводах допускаются общие корни \(D,W\).

Восстановление по концевой спектральной мере и по известной части матрицы Якоби относится к классической обратной спектральной задаче Хохштадта; см. Гестези и Саймон \cite[теорема 4.2]{gs}. Формула \eqref{eq:split} даёт прямую специализацию для соседних измерительных точек. Малый случай \(n=5\) дополнительно использует амплитуды наблюдений.

\newpage
\section{Произвольный выход с хвостом длины не менее двух}
В этом разделе спектральная переменная \(x\) относится к \(xI-T\).
Положим \(F(x)=(-1)^nD(-x)\), \(G(x)=(-1)^{n-3}W(-x)\);
через \(B_j(x)=(-1)^jL_j(-x)\) обозначим левый полином. При чтении матрицы от конца положим \(P_d(x)=\det(xI-T_{[n-d+1,n]})\), так что \(P_{n-3}=G\). Пусть \(x_i\) --- простые собственные значения \(T\), а \(\rho_i>0\) --- нормированные массы спектральной меры в вершине \(n\). По трёхчленной рекурсии \(P_d(T)e_n=C e_{n-d}\) для \(d<n\), где \(C=\prod_{v=n-d}^{n-1}(-\sqrt{\beta_v})\). В частности \(P_d\) ортогонален всем полиномам степени меньше \(d\).

Предположим, что \(F,G\) взаимно просты. Для всех реализаций с заданными \(F,G\)
\begin{equation}\label{eq:rho}
 \rho_i=\frac{U(x_i)}{F'(x_i)G(x_i)},\qquad U=K B_2,\qquad
 K=\prod_{v=3}^{n-1}\beta_v>0 .
\end{equation}
Это формула вычетов из \cite[§3]{v2}: тождество для дополнительных миноров даёт \(G B_{n-1}-F\det(xI-T_{[4,n-1]})=K B_2\), а \(B_{n-1}/F\) есть концевая резольвента. Нормированные массы однозначно определяют \(T\) по ортогонализации \(1,x,\ldots,x^{n-1}\) при заданных знаках внедиагональных элементов.

Пусть \(d=n-r\ge2\), и две реализации дают те же наблюдения. Они имеют общий \(P_d\). Если \(c_1,c_2\) --- корни \(B_2\), то \(c_1\ne c_2\) по \(\beta_1>0\); из \eqref{eq:rho} ни один \(c_j\) не есть узел \(x_i\). Для числителя \(V\) второй реализации имеем
\[
 \frac{V(x)}{U(x)}=a+\frac{b}{x-c_1}+\frac{c}{x-c_2},\qquad
 I_j=\sum_i\frac{\rho_iP_d(x_i)}{x_i-c_j}.
\]
Ортогональность \(P_d\) к \(1,x\) для обеих мер даёт \(bI_1+cI_2=0\) и \(c_1 bI_1+c_2 cI_2=0\). Если \(I_1I_2\ne0\), то \(b=c=0\); нормировка даёт \(a=1\). Обе меры и все скорости совпадают.

Кофактор резольвенты и равенство \(P_d(T)e_n=C e_r\) дают
\[
 \sum_i\frac{\rho_iP_d(x_i)}{z-x_i}
   =C^2\frac{B_{r-1}(z)}{F(z)}.
\]
Таким образом, достаточно взаимной простоты \(B_2,B_{r-1}\). Для \(r=1,2,4\) она автоматическая. При \(r\ge5\) в корне \(c_j\) полинома \(B_2\) разложение по ребру 2 даёт \(B_{r-1}(c_j)=-\beta_2B_1(c_j)\det(c_jI-T_{[4,r-1]})\). Соответствующий результант не тождественно равен нулю: можно расположить спектры непересекающихся блоков \([1,2]\) и \([4,r-1]\) в непересекающихся интервалах, выбирая разные положительные диагонали и малые ненулевые связи. Произвольная такая положительно определённая \(T\) физически достижима при любой \(\ell\): при \(w=T^{-1}e_\ell>0\) матрица \(-\diag(w)T\diag(w)^{-1}\) имеет нужные положительные скорости и единственную утечку \cite[§2]{v2}.

\newpage
\section{Что определяет измерение в предпоследней камере}
Пусть \(r=n-1\), \(n\ge6\). Отношение \(G_{n-1}/G_n=(s+d_n)/a_{n-1}\) даёт \(d_n,a_{n-1}\). При \(\ell=n\) по балансу \(h=D(0)/q\); для остальных \(\ell\) в камере \(n\) нет выведения. В обоих случаях известны \(b_{n-1}=d_n-h\mathbf1_{\ell=n}\) и \(\beta_{n-1}\). Значит, нам известны три концевых момента
\[
 \sum_i \rho_i=1,\qquad \sum_i x_i\rho_i=d_n,\qquad
 \sum_i x_i^2\rho_i=d_n^2+\beta_{n-1}.
\]
При \(n=6\) к ним добавляются три нулевых момента полиномов \(G,xG,x^2G\) по ортогональности. Вместе с \(1,x,x^2\) это базис полиномов степени до пяти, поскольку \(G\) монический кубический. При шести различных известных узлах значения этих шести функционалов однозначно определяют веса по Вандермонду. Это доказательство допускает и общие корни \(F,G\).

В общем случае на взаимно простых \(F,G\) три момента образуют линейную систему для трёх коэффициентов \(U\) из \eqref{eq:rho}. Её матрица равна \((M_{j+k})_{j,k=0}^2\), где
\[
 M_v=\sum_i\frac{x_i^v}{F'(x_i)G(x_i)}.
\]
Покажем, что её определитель типично отличен от нуля. Эти величины рациональны по коэффициентам \(F,G\), например \(M_v=[x^{n-1}](x^v G^{-1}\bmod F)\) по формуле Лагранжа. Для \(v\le4\) сумма вычетов \(x^v/(FG)\) в конечных полюсах равна нулю. Взяв \(G\) с \(m=n-3\ge3\) разными вещественными корнями \(w\) и \(F(w)=-G'(w)\), получим
\[
 M_v=\sum_{G(w)=0}\frac{w^v}{G'(w)^2}.
\]
Такие монические \(F\) существуют: можно взять \(F=GJ-G'\) с монической кубикой \(J\); добавлением константы к \(J\) избегаем кратных корней \(F\), так как они возможны лишь в критических точках рациональной функции \(J-G'/G\). Последняя матрица моментов положительно определена на полиномах степени до двух. Следовательно, числитель определителя не тождественно нулевой.

Чтобы перенести вывод на физические модели, существенно, что достижимые пары \(F,G\) имеют внутренность в пространстве монических коэффициентов. Например, выберем положительные простые узлы \(F\) и поместим по одному корню \(G\) в произвольные \(n-3\) разные промежутки. Поместим два корня монической квадратичной \(Q\) в оставшиеся промежутки. Тогда \(\rho_i\propto Q(x_i)/(F'(x_i)G(x_i))>0\); ортогонализация даёт \(P_{n-3}=G\), так как \(\sum x_i^jQ(x_i)/F'(x_i)=0\) при \(j<n-3\). Строгие неравенства сохраняются на открытом множестве. Полученный \(T\) положительно определён и допускает физическую реализацию при каждой \(\ell\). Это заодно доказывает типичную взаимную простоту \(F,G\), использованную в §3. Коэффициенты миноров полиномиальны по \(d_j,\beta_j\), а значит, и по исходным скоростям. После исключения нулей указанных ненулевых полиномов восстановление единственно среди всех положительных реализаций.

\newpage
\section{Построение пар, неразличимых по трём выходам}
Докажем необходимость условия для гарантии, равномерной по месту выведения. Здесь \(\ell=1\), \(n\ge7\), \(5\le r<n\) и \(d=n-r\). Фиксируем положительно определённую исходную \(T\) с взаимно простыми \(F,G\). На нормированной плоскости числителей \(U\) из \eqref{eq:rho} фиксируем старший коэффициент \(K\). Если нормировка и \(K\) независимы, положительные веса образуют открытый ограниченный отрезок: они меняются аффинно, положительны в исходной точке и имеют сумму один. Координаты различны по интерполяции \(U\) на полном спектре. Следующие два факта позволяют построить нужную пару на этом отрезке.

\smallskip\noindent\textit{1. Равенство амплитуд.}
Обозначим \(\|P_d\|^2=\sum_i\rho_iP_d(x_i)^2=\prod_{v=r}^{n-1}\beta_v\), где последняя формула следует из \(P_d(T)e_n=C e_r\). Для \(\ell=1\) рекурсия справа при \(s=0\) даёт \(R_{r+1}(0)=\prod_{v=r}^{n-1}b_v\). Поэтому
\[
    g_r=q\,R_{r+1}(0)/\|P_d\|^2,\qquad p=qW(0)/K.
\]
Если \(c\) --- корень \(B_2\) и \(B_{r-1}(c)=0\), то \(\sum_i\rho_iP_d(x_i)/(x_i-c)=0\) по кофактору из §3. Замена \(\rho_i\) на \(\rho_i(1+t/(x_i-c))\) сохраняет \(P_d\) и его норму, пока веса положительны и нормированы. Для ортогональности разделим \(x^j\) на \(x-c\), \(j<d\); полиномиальная часть ортогональна \(P_d\). Для нормы таким же образом разделим \(P_d(x)-P_d(c)\) на \(x-c\). Аргумент линеен и допускает сумму двух таких возмущений.

\smallskip\noindent\textit{2. Два регулярных решения.}
При \(\ell=1\) на отрезке весов справедливо \cite[лемма 3]{v2}
\[
 \frac{q}{D(0)}=\min_{\deg f\le n-1,\ f(0)=1}
                          \sum_i\frac{\rho_i}{x_i} f(x_i)^2 .
\]
Внутри отрезка минимум положителен: полином допустимой степени не исчезает в \(n\) узлах. На каждом конце хотя бы одна масса обращается в нуль. Полином с корнями в оставшихся узлах и значением 1 в нуле показывает, что минимум стремится к нулю. Минимум вогнут и дифференцируем внутри (положительно определённая матрица Грама). На любом уровне строго между нулём и максимумом есть два пересечения с ненулевой производной: знак производной в каждом следует из наклона секущей к точке максимума. Оба набора имеют одно и то же \(p,q\). Они регулярны для карты \((D,W,p,q)\): при взаимной простоте коэффициенты \(D,W\) и два свободных коэффициента нормированного \(U\) служат локальными координатами \cite[лемма 2]{v2}; на плоскости независимы \(\dd K,\dd q\), а \(p=qW(0)/K\).

\smallskip\noindent\textit{Конструкция при \(5\le r\le n-2\).}
Выберем у левого двухкамерного блока собственные значения \(c=3,e=5\) (диагональ \(4,4\), связь \(-1\)). У обоих непустых блоков \([4,r-1]\), \([r+1,n-1]\) зададим матрицу \(3I+\tfrac18\mathcal L\), где \(\mathcal L\) --- лапласиан пути; у одноточечного блока это просто \(3\). Остальные диагонали равны 3, остальные связи \(-1/8\). Матрица положительно определена по строгому диагональному преобладанию; весь блок \([4,n]\) имеет спектр ниже 5. У \([4,n-1]\) есть собственное значение 3: на двух блоках берём противоположные константы, в вершине \(r\) --- нуль. По строгому чередованию с соседним хвостом у \([4,n]\) этого значения нет.

\newpage
Разложение по ребру 2 при \(z=c,e\) даёт
\(B_j(z)=-\beta_2 B_1(z)\det(zI-T_{[4,j]})\) для \(j\ge3\).
Значит, \(B_{r-1}(c)=B_{n-1}(c)=0\), \(B_{n-1}(e)\ne0\); \(F,G\) взаимно просты, поскольку их общий корень обязан быть корнем \(B_2\) по тождеству миноров, а \(G(c)G(e)\ne0\). При этом \(F(c)F(e)\ne0\) по указанному разложению. Из \(B_{n-1}/F=\sum_i\rho_i/(x-x_i)\) следует, что направление \(U/(x-c)\) сохраняет нормировку, а \(U/(x-e)\) меняет её. Оба имеют степень один. Следовательно, нормировка независима от старшего коэффициента, а отрезок при фиксированном \(K\) параметризуется деформацией из пункта 1 с полюсом \(c\). Оба пересечения из пункта 2 имеют также общий \(G_r\).

\smallskip\noindent\textit{Конструкция при \(r=n-1\).}
Возьмём неприводимый симметричный трёхдиагональный блок на \([4,n-2]\), имеющий размер не меньше двух, и выберем два его различных собственных значения \(c,e\). У левого двухкамерного блока зададим те же значения, например одинаковые диагонали \((c+e)/2\) и связь \(-|c-e|/2\). Достроим цепочку с ненулевыми отрицательными связями и \(d_n\notin\{c,e\}\), затем сдвинем всю диагональ на одну константу до положительной определённости, сдвигая \(c,e\) вместе с ней. По чередованию у \([4,n-1]\) нет \(c,e\), а по последней трёхчленной рекурсии их нет и у \([4,n]\). Таким образом, \(F,G\) взаимно просты, \(B_{r-1}(c)=B_{r-1}(e)=0\), \(B_{n-1}(c)B_{n-1}(e)\ne0\). Нормировка снова независима от \(K\). На построенном отрезке при фиксированных \(F,G,K\) отношение нового числителя к исходному имеет вид \(1+b/(x-c)+b'/(x-e)\). По пунктам 1 и 2 получаем пару со всеми тремя совпадающими выходами. Это завершает доказательство теоремы~\ref{t:main}.

\section{Семикамерный рациональный контрпример}
Отдельный пример показывает, что даже скорость выведения может различаться. Пусть \(n=7,\ \ell=5,\ r=6\). В положительной спектральной переменной зададим
\[
 G=(x-3)(x-7)(x-9)(x-16),\quad V=(x-4)(x-10)(x-17),\quad
 F=VG+80(x-5)(x-19).
\]
Пусть \(J=xG^{-1}\bmod F,\ H_0=J/[x^6]J\), где \([x^6]J=287/10137600\).
Для двух значений \(u=9000,\ 971185627456000/106036807641\) положим
\(H=H_0-uV/80\). Это задание пары рациональными полиномами: по \(F,H\) деление соседних миноров однозначно восстанавливает обе цепочки, как в §2. Прямой пересчёт коэффициентов даёт в обоих случаях
\[
 p=95,\quad q=\frac{2674979300}{82101},\quad
 a_6=\frac{1504800}{531811},\quad b_6=\frac{1853}{287}.
\]
Скорости выведения различаются: в указанном порядке они равны
\[
  \frac{38362500957359744}{5672745222405219},\qquad
  \frac{149303720272872436416340353}{21445474271243680437492328}.
\]
Положительность скоростей, равенство \(G_3,G_6,G_7\) и знаки двух якобианов порядка 13 для исходной карты \((D,W,p,q)\), соответственно \(-,+\), проверяет \texttt{verify\_seven\_collision.py}. В коде единичная конечная разность даёт точную производную: каждая скорость входит лишь в один столбец \(sI-A\), а каждое произведение \(p,q\) аффинно по каждой скорости в отдельности.

\newpage
\section{Слепое восстановление по двум выходам}
Для взаимно простых \(D,W\) с рациональными коэффициентами программа получает только \(n,\ell,D,W,p,q\). Следующий алгоритм перечисляет все положительные регулярные решения. Используем переменную Лапласа \(s\), \(P(s)=c_0+c_1s+c_2s^2\), \(H=(P W^{-1}\bmod D)\). По формуле миноров \(H=L_{n-1}\). Моничность \(H\) даёт линейное уравнение; для \(\ell\ge3\) добавляется \(c_0=pc_2\), для \(\ell\le2\) --- \(c_2=qW(0)/p\). На непустом регулярном слое остаётся аффинная прямая. Моменты \(\mu_j=\sum\rho_i x_i^j\) извлекаются из \(H/D\) в бесконечности и \(\mu_{-1}=H(0)/D(0)\).

При \(\ell=3\) оставшееся условие есть \(c_2=qW(0)/p\), при \(\ell=n\) --- \(H(0)=q\). В остальных случаях \(m=n-\ell\), \(b=(\mu_j)_{j=0}^{m-1}\), \(C=(\mu_{i+j+1})_{i,j=0}^{m-1}\), и решается уравнение
\[
 \det\begin{pmatrix}\mu_{-1}-q/D(0)&b^\top\\ b&C\end{pmatrix}=0 .
\]
На положительных реализациях \(C>0\); дополнение Шура воспроизводит вариационный минимум \cite[§4]{v2}. На прямой простой корень эквивалентен ненулевому якобиану в локальных координатах. Нулевое уравнение или зависимые согласованные линейные условия дают только критические прообразы. Для каждого действительного корня программа делением \(D,H\) восстанавливает матрицу, проверяет положительность скоростей и повторно вычисляет хвост и обе амплитуды. Изоляция и знаки вычисляются в точной рациональной или вещественной алгебраической арифметике. Общий множитель \(D,W\) помечается отдельным статусом.

В тесте два набора данных из \cite[§§6, 8]{v2} передаются программе по очереди. Пусть \(\Delta N_4=D(G_4^{(1)}-G_4^{(2)})\), порядок реализаций как в исходной работе.
\begin{center}\small
\begin{tabular}{llll}
\toprule
\(n,\ell\) & \(W;\ (p,q)\) & якобианы двух решений & \(\Delta N_4\)\\
\midrule
\(5,4\) & \(s^2+7s+8;\ (2,6)\) & \(3456,\ -4698\) & \(3(s-1)/2\)\\
\(5,4\) & \(s^2+11s+6;\ (12,192)\) & \(-2359296,\ 3760128\) & \(8s\)\\
\bottomrule
\end{tabular}
\end{center}
Обе пары различает и каждая из камер 1 и 2. Во втором примере обе скорости выведения равны 2 и значения \(G_4(0)\) совпадают. Измерение одной стационарной величины было бы недостаточным.

\smallskip
Классификация исключений отдельно для каждого фиксированного \(\ell\) и устойчивость восстановления по шумным пробам остаются отдельными вопросами.

\section*{Сопоставление с предшествующими результатами}
Наиболее прямой предшественник §3 --- уже опубликованный результат автора \cite{endp}. Для входа в камере 1, двух выходов в \(j<k\) и любой заданной единственной утечки условие \(j-1\le n-k\) достаточно для типичной глобальной единственности. Беря \(j=3,k=r\), получаем ровно \(4\le r\le n-2\); измерение в камере \(n\) в этом выводе не требуется. Для доказательства типичной единственности при \(4\le r\le n-2\) есть также вывод из ранее опубликованных результатов. В \cite[лемма 2.1, §6]{cg} дополнительные вложенные спектры дают линейные условия ортогональности на нормированные концевые веса. Для \(o\in\{3,r,n\}\) эти условия имеют строки \((x_i^j P_{n-o}(x_i))_i,\ 0\le j<n-o\). Тождество для миноров в узлах \(F\) даёт
\[
 B_{n-1}(x_i)P_{n-o}(x_i)=K_o B_{o-1}(x_i),
 \qquad K_o=\prod_{v=o}^{n-1}\beta_v .
\]
Концевые миноры \(B_{n-1}\) и \(F\) взаимно просты. Умножение столбцов на \(B_{n-1}(x_i)\) и деление строк на \(K_o\) сводят задачу о ранге к вычислению размерности пространства
\[
 B_2\,\mathcal P_{\le n-4}+B_{r-1}\,\mathcal P_{\le n-r-1},
\]
где \(\mathcal P_{\le t}\) --- полиномы степени до \(t\). Все его элементы имеют степень не выше \(n-2\), поэтому вычисление в \(n\) различных узлах инъективно. Для \(4\le r\le n-2\) пространство кратных каждого из двух полиномов здесь заполняет всю допустимую по степени часть его главного идеала. Пересечение состоит из кратных НОК; формула сложения размерностей даёт \(n-1-\deg\gcd(B_2,B_{r-1})\). Типичный ранг \(n-1\) следует также из односторонней рекурсии \cite[приложение C.2, с.~22]{anon}. При достижимых данных нормировка независимо добавляет одну строку, поскольку положительные массы имеют сумму 1. Получаем единственные веса и единственные скорости.

Таким образом, для этого диапазона типичная глобальная единственность является следствием сочетания двух прежних результатов с тождеством миноров. Теорема~1 самой работы \cite{anon} относится к локальной идентифицируемости. При \(r=n-1\) дополнительная амплитуда в §4 даёт второй момент; приведённый выше расчёт ранга на этот случай не переносится. У Michor и Teschl \cite[теорема 2.1]{mt} измеряются спектры двух блоков после удаления одной вершины. Ни один из этих выводов сам по себе не даёт конструкции регулярных положительных пар §5.

\smallskip\noindent\textit{Дополнительная проверка регулярности из §5.}
Пусть \(\lambda_i^L,\rho_i\) --- массы одного и того же узла \(x_i\) на левом и правом концах, \(C_*=\prod_{j<n}\beta_j\). Концевой кофактор \((1,n)\) резольвенты имеет постоянный числитель с квадратом \(C_*\). Поэтому вычет в простом собственном значении (произведение координат нормированного собственного вектора) даёт
\(\lambda_i^L\rho_i=C_*/F'(x_i)^2\). При \(\ell=1\) имеем \(q=C_*/R_2(0)\) и \(R_2(0)/D(0)=\sum_i\lambda_i^L/x_i\). Следовательно,
\[
        \frac{D(0)}q
          = D(0)^2\sum_i\frac{1}{x_i F'(x_i)^2\rho_i}.
\]
На построенном невырожденном аффинном отрезке положительных весов правая часть строго выпукла: хотя бы один вес имеет ненулевую производную. Она стремится к \(+\infty\) на обоих концах. Каждый уровень выше минимума даёт две точки с противоположными ненулевыми производными. Это даёт второе доказательство трансверсальности.

\begin{thebibliography}{9}\small
\bibitem{v2} M. Aharodni. \href{https://github.com/qqhgdkktu/catenary-identifiability/releases/tag/global-multiplicity-v2}{Positive regular fibers of a catenary compartmental model with two outputs}. Препринт, версия 2 (2026).
\bibitem{endp} M. Aharodni. \textit{Two outputs in a catenary compartmental model}. \href{https://github.com/qqhgdkktu/catenary-identifiability/blob/7300631e25bea9a7afe03e6426122d93010a423d/Chain/EndpointExterior.lean\#L329-L341}{теорема \texttt{generic\_global\_small}}, 24 ноября 2026 г., версия \texttt{7300631}.
\bibitem{anon} Anonymous. \href{https://doi.org/10.5281/zenodo.22909460}{A placement criterion for catenary compartment models: Generic local identifiability on a labelled bidirected path}. Препринт, 23 сентября 2026 г.
\bibitem{gs} F. Gesztesy, B. Simon. \href{https://doi.org/10.1007/BF02788147}{\(m\)-functions and inverse spectral analysis for finite and semi-infinite Jacobi matrices}. J. Analyse Math. 73 (1997), 267--297.
\bibitem{cg} K. Castillo, G. Gordillo-Núñez. \href{https://arxiv.org/abs/2602.20894}{Wendroff's theorem beyond consecutive degrees and related inverse spectral problems}. arXiv:2602.20894 (2026).
\bibitem{mt} J. Michor, G. Teschl. \href{https://arxiv.org/abs/math/0211051}{Reconstructing Jacobi matrices from three spectra}. Oper. Theory Adv. Appl. 154 (2004), 151--154.
\end{thebibliography}
\end{document}
